% !TEX program = xelatex
% ============================================================================
% Propuesta delimitada de dimensiones del sistema
% Proyecto Final de Bases de Datos - Inventario, Compras y Ventas (PostgreSQL)
% ============================================================================
\documentclass[11pt,oneside]{article}

\usepackage{estilo-proyecto}
\docnombre{Dimensiones del sistema}

\begin{document}

\portadilla{Proyecto final \textperiodcentered{} Bases de datos}
  {Dimensiones del sistema}
  {Propuesta delimitada de datos \textperiodcentered{} PostgreSQL}

\section{Descripci\'on}

Este documento fija las dimensiones de datos que guiar\'an la poblaci\'on inicial
de la base de datos y las pruebas funcionales del sistema. Para cada variable se
definen dos valores:

\begin{itemize}[leftmargin=1.5em, itemsep=2pt]
  \item \textbf{Requerimiento base}: m\'inimo necesario para demostrar las
        funcionalidades del sistema.
  \item \textbf{Propuesta delimitada}: cantidad realista que permite pruebas de
        carga leve, consultas por rango de fechas y verificaci\'on de las
        particiones LIST.
\end{itemize}

La propuesta delimitada se adopta como est\'andar para la generaci\'on de datos
de prueba.

\newpage
\section{Tabla de dimensiones}

\renewcommand{\arraystretch}{1.4}
\begin{tabular}{%
    >{\sffamily\small\raggedright\arraybackslash}p{5.5cm}
    >{\centering\small\arraybackslash}p{5.1cm}
    >{\centering\small\bfseries\arraybackslash}p{5.1cm}}
\toprule
\rowcolor{fondo}
{\sffamily\bfseries\footnotesize Variable} &
{\sffamily\bfseries\footnotesize Requerimiento base} &
{\sffamily\bfseries\footnotesize Propuesta delimitada} \\
\midrule
Sedes totales                      & 4 (1 de distribuci\'on) & 5 (1 distribuci\'on, 4 ventas) \\
Categor\'ias de producto           & 5                       & 6 \\
Productos activos                  & 80                      & 100 \\
Clientes registrados               & 100                     & 120 \\
Empleados y usuarios               & 20                      & 25 \\
Proveedores                        & 10                      & 12 \\
Almacenes y bodegas                & 6                       & 8 \\
Pedidos de venta                   & 300                     & 350 \\
\'Ordenes de compra                & 60                      & 75 \\
Devoluciones                       & 20                      & 25 \\
Transferencias entre almacenes     & 25                      & 30 \\
Registros de auditor\'ia           & 500                     & 1\,500 \\
Rango hist\'orico transaccional    & 4 meses                 & 4 meses \\
\bottomrule
\end{tabular}

\section{Notas de implementaci\'on}

\begin{enumerate}[leftmargin=1.5em, itemsep=3pt]
  \item \textbf{Particiones LIST.} Las tablas \texttt{pedidos\_venta} (esquema
        \texttt{ventas\_pagos}) e \texttt{inventarios} (esquema
        \texttt{inventarios\_\allowbreak transferencias}) contemplan fragmentaci\'on por
        \texttt{id\_sede}. Con cinco sedes se definen cinco fragmentos por tabla,
        lo que permite validar consultas por sede y vistas con \texttt{UNION ALL}.

  \item \textbf{Cobertura de roles.} Los veinticinco empleados y usuarios cubren los
        roles definidos en el sistema: administrador, gerencia, inventario, compras,
        vendedor y auditor.

  \item \textbf{Relaci\'on \textit{suministra}.} Con doce proveedores y cien
        productos se puebla la tabla \texttt{producto\_proveedor} (en
        \texttt{catalogo\_terceros}) con cien registros iniciales, uno por
        producto y abierto a m\'ultiples proveedores por referencia.

  \item \textbf{Consistencia de existencias.} Los 350 pedidos de venta y las 75
        \'ordenes de compra se distribuyen en los cuatro meses del hist\'orico
        para que las consultas de alertas de reabastecimiento y la vista de
        ventas mensuales devuelvan datos en todos los per\'iodos.

  \item \textbf{Integridad referencial.} Todas las devoluciones referencian
        pedidos existentes y completados, y todas las transferencias enlazan
        almacenes distintos (\texttt{id\_almacen\_origen <> id\_almacen\_destino}).
\end{enumerate}

\end{document}
